%%%%%%%%%%%%%%%%%%%%%%%%%%%%%%%%%%%%%%%%%%%%%%%%%%%%%%
% A Beamer template for University of Wollongong     %
% Based on THU beamer theme                          %
% Author: Qiuyu Lu                                   %
% Date: July 2024                                    %
% LPPL Licensed.                                     %
%%%%%%%%%%%%%%%%%%%%%%%%%%%%%%%%%%%%%%%%%%%%%%%%%%%%%%
% Customized for Sharif University of Technology     %
%%%%%%%%%%%%%%%%%%%%%%%%%%%%%%%%%%%%%%%%%%%%%%%%%%%%%%


\documentclass[serif, aspectratio=169]{beamer}
%\documentclass[serif]{beamer}  % for 4:3 ratio
\usepackage[T1]{fontenc} 
\usepackage{fourier} % see "http://faq.ktug.org/wiki/uploads/MathFonts.pdf" for other options
\usepackage{hyperref}
\usepackage{latexsym,amsmath,xcolor,multicol,booktabs,calligra}
\usepackage{graphicx,pstricks,listings,stackengine}
\usepackage{lipsum}
\usepackage{tikz}
% For writing clean pseudocodes
% \usepackage{algorithm, algpseudocode, mathtools, needspace}
% \usepackage{algorithmic}
\usepackage{algorithm}
\usepackage{algorithmic}
\usepackage{algpseudocode}

% To justify the items


\author{Ali Sharifi-Zarchi}
% \author{CE Department}
\title{Machine Learning (CE 40477)}
\subtitle{Fall 2024}
\institute{
    CE Department \\
    Sharif University of Technology
}
%\date{\small \today}
% \usepackage{UoWstyle}
\usepackage{SUTstyle}

% defs
\def\cmd#1{\texttt{\color{red}\footnotesize $\backslash$#1}}
\def\env#1{\texttt{\color{blue}\footnotesize #1}}
\definecolor{deepblue}{rgb}{0,0,0.5}
\definecolor{deepred}{RGB}{153,0,0}
\definecolor{deepgreen}{rgb}{0,0.5,0}
\definecolor{halfgray}{gray}{0.55}

\lstset{
    basicstyle=\ttfamily\small,
    keywordstyle=\bfseries\color{deepblue},
    emphstyle=\ttfamily\color{deepred},    % Custom highlighting style
    stringstyle=\color{deepgreen},
    numbers=left,
    numberstyle=\small\color{halfgray},
    rulesepcolor=\color{red!20!green!20!blue!20},
    frame=shadowbox,
}

% For writing comments that are aligned to the left side
\makeatletter
\NewDocumentCommand{\LeftComment}{s m}{%
  \Statex \IfBooleanF{#1}{\hspace*{\ALG@thistlm}}\(\triangleright\) #2}
\makeatother
% To manually indent states in algorithmicx
\newcommand{\IndState}{\State\hspace{\algorithmicindent}}
% To make breakable algorithms
\makeatletter
\newenvironment{nofloatalgorithmic}[2][0]
  {
  \par
  \needspace{\dimexpr\baselineskip+6.8pt}
  \noindent
  \hrule height.8pt depth0pt \kern2pt
  \refstepcounter{algorithm}
  \addcontentsline{loa}{algorithm}{\numberline{\thealgorithm}#2}
  \noindent\textbf{\fname@algorithm~\thealgorithm} #2\par
  \kern2pt\hrule\kern2pt
  \begin{algorithmic}[#1]
  }
  {
  \end{algorithmic}
  \nobreak\kern2pt\hrule\relax
  }
\makeatother
% To make vertical arrow
\newcommand\vertarrowbox[3][6ex]{%
  \begin{array}[t]{@{}c@{}} #2 \\
  \left\uparrow\vcenter{\hrule height #1}\right.\kern-\nulldelimiterspace\\
  \makebox[0pt]{\scriptsize#3}
  \end{array}%
}
% Clean argmin
\DeclareMathOperator*{\argmin}{arg\,min}


\begin{document}

\begin{frame}
    \titlepage
    \vspace*{-0.6cm}
    \begin{figure}[htpb]
        \begin{center}
            \includegraphics[keepaspectratio, scale=0.25]{pic/sharif-main-logo.png}
        \end{center}
    \end{figure}
\end{frame}

\begin{frame}    
\tableofcontents[sectionstyle=show,
subsectionstyle=show/shaded/hide,
subsubsectionstyle=show/shaded/hide]
\end{frame}

\section{Unsupervised Learning Overview}
\begin{frame}{Unsupervised Learning}
    \textbf{Unsupervised Learning} involves working with \textbf{unlabeled data}, where the goal is to \textbf{infer the natural structure} present within a set of data points.
    \begin{itemize}
        \item Learning from unlabeled data.
        \item Most of the times, there is no (or minimal) prior knowledge of the data.
        \item Two of the most common techniques:
        \begin{itemize}
            \item \textbf{Clustering}: Grouping data points into clusters based on similarity towards user need.
            \item \textbf{Dimensionality Reduction}: Reducing the number of features under consideration and keeping (perhaps approximately) the most informative features.
        \end{itemize}
    \end{itemize}
\end{frame}

\begin{frame}{Clustering: Bio-informatics}
    \begin{figure}
        \centering
        \includegraphics[scale=0.5]{pic/bioconductor_clustering.png}
    \end{figure}
        \begin{tikzpicture}[remember picture,overlay]
        \node[anchor=south west, xshift=0.1cm, yshift=0.22cm] at (current page.south west) {
            \scriptsize Figure adapted from \href{https://bioconductor.org/books/3.15/OSCA.basic/clustering.html}{bioconductor.org}
        };
    \end{tikzpicture}
\end{frame}

\begin{frame}{Slide showcasing reduction (PCA)}
    
\end{frame}

\begin{frame}{Clustering: Customer segmentation}
    \begin{figure}
        \centering
        \includegraphics[scale=0.45]{pic/customer_clusters_raw_plot.png}
    \end{figure}
    \begin{center}
        Information of customers
    \end{center}
\end{frame}

\begin{frame}{Clustering: Customer segmentation (cont.)}
    \begin{figure}
        \centering
        \includegraphics[scale=0.45]{pic/customer_clusters.png}
    \end{figure}
    \begin{center}
        Predicted clusters of customers
    \end{center}
\end{frame}


\section{K-Means}
\begin{frame}{K-Means overview}
    \begin{itemize}
        \item One of the most common partitional clustering methods used
        \item The idea is to \textbf{find} \( \mathbf{K} \) \textbf{centers.} \textbf{Each center representing a cluster.}
        \item Each data point is assigned to cluster \( j \) if and only if it has the least distance to center of cluster \( j \) amongst all clusters.
        \item K-Means suggest an \textbf{iterative algorithm} to find these centers.
    \end{itemize}
\end{frame}

\begin{frame}{Problem Intuition}
    \begin{itemize}
        \item One of the most straightforward tasks we can perform on a data set without labels is to find groups of data in our dataset which are "similar" to one another -- what we call clusters.
        \item  In a space, where each data point can be associated with a vector $x_i \in \mathbb{R}^d$. Lets formalize and model this problem. What is a cluster ?
        \item Intuitively (and naively), we might think of a cluster as a subset of data points whose inter-point euclidean distances are small compared with the distances to points outside of the cluster.
        \item For now lets fix a $k$, meaning that we want to find $k$ clusters, or we know there are $k$ clusters.
        \item Also we might want to be able to predict the cluster a new data-point belongs to.
        
    \end{itemize}
\end{frame}


\begin{frame}{Problem definition}
    \begin{itemize}
        \item Like we did before in supervised learning, we want to create an objective function or loss function. But here our objective is assessing the clusters we have created. Formally:
        \item We have $X_{\text{train}} = \{ x^{(1)}, x^{(2)}, \dots, x^{(N)} \} \subseteq \mathbb{R}^d$
        \item  We are learning:
        \begin{enumerate}
            \item a function $f:\mathbb{R}^d\to \{1,2, \dots , K\}$ that assigns a cluster to each data point. 
\item  a set of \( K \) prototypes \(\mu = \{ \mu_1, \mu_2, \dots, \mu_K  \} \subseteq \mathbb{R}^d \) as the cluster  representatives.
        \end{enumerate}
        \item For now we lets create the objective as below:
        $$
        \sum_{\mathbf{x} \in X_{\text{train}}} || x - \mu_{f(x)} ||^2
        $$
        
    \end{itemize}
\end{frame}


\begin{frame}{Problem definition (cont.)}
    \begin{itemize}
        

        \item  We can express $f$ by defining $r_k(\mathbf{x}) = 1$ if $f(\mathbf{x}) = k$ and $0$ otherwise, we can write this objective as below:
$$
J = \sum_{\mathbf{x} \in X_{\text{train}}} \sum_{k=1}^{K} r_k(x)|| x -  \mu_k ||^2
$$
        \item This is also known as \textit{distortion measure}
        \item We want to chose $f$ and $\mathbf{C}_k$ to minimize this measure.
        \item K-Means is an algorithm to do this task.

        
    \end{itemize}
\end{frame}


\begin{frame}{K-Means }
    \begin{itemize}
   

\item If we fix a set of centroids or representetives $\mu$ then we could optimize by \textbf{assigning} $f(x) : = \textit{argmin}_k \; ||x - \mu_k||^2$ and $r_k(x)$ with respect to this $f$.
\item If we fix an assignment $f$ then we can optimize by taking derivative and \textbf{updating} $\mu_k = \frac{\sum_i r_k(x_i) x_i}{\sum_i r_k(x_i) }$.
\item K-Means uses an iterative process that:
\begin{enumerate}
    \item \textbf{Assigns} each point to the \textbf{nearest} centroid (cluster representative).
    \item \textbf{Updates} each centroid as the \textbf{mean} of the points assigned to that cluster.

\end{enumerate}
    \item GIF Here
    \end{itemize}
\end{frame}

\begin{frame}{K-Means Algorithm}
    \begin{itemize}
        The K-means algorithm can be expressed as:
    \end{itemize}
    
\begin{algorithm}
\caption{K-means Clustering Algorithm}
\begin{algorithmic}[1]
\State \textbf{Input:} Data points $X = \{x_1, x_2, ..., x_n\} \subset \mathbb{R}^d$, number of clusters $K$
\State \textbf{Output:} Cluster assignments $C = \{C_1, C_2, ..., C_K\}$ and centroids $\mu = \{\mu_1, \mu_2, ..., \mu_K\}$
\State \textbf{Initialize:} Randomly select $K$ points as initial centroids $\mu_1, \mu_2, ..., \mu_K$
\Repeat
    \State \textbf{Assignment step:}
    \For {each data point $x_i$}
        \State Assign $x_i$ to the nearest centroid $C_k$: 
        \[
        C_k = \{ x_i : ||x_i - \mu_k||^2 \leq ||x_i - \mu_j||^2, \, \forall j = 1, 2, ..., K \}
        \]
    \EndFor
    \State \textbf{Update step:}
    \For {each cluster $C_k$}
        \State Update the centroid of $C_k$: 
        \[
        \mu_k = \frac{1}{|C_k|} \sum_{x_i \in C_k} x_i
        \]
    \EndFor
\Until{the centroids do not change (convergence)}
\end{algorithmic}
\end{algorithm}
\end{frame}

% K means in action
\begin{frame}{K-Means in action}
    \begin{figure}
        \centering
        \includegraphics[scale=0.45]{pic/figs/kmeans_step_1_unlabeled_data.png}
    \end{figure}
\end{frame}
\begin{frame}{K-Means in action (cont.)}
    \begin{figure}
        \centering
        \includegraphics[scale=0.45]{pic/figs/kmeans_step_2_assignment.png}
    \end{figure}
\end{frame}

\begin{frame}{K-Means in action (cont.)}
    \begin{figure}
        \centering
        \includegraphics[scale=0.45]{pic/figs/kmeans_step_3_recompute_centers.png}
    \end{figure}
\end{frame}
\begin{frame}{K-Means in action (cont.)}
    \begin{figure}
        \centering
        \includegraphics[scale=0.45]{pic/figs/kmeans_step_4_assignment.png}
    \end{figure}
\end{frame}
\begin{frame}{K-Means in action (cont.)}
    \begin{figure}
        \centering
        \includegraphics[scale=0.45]{pic/figs/kmeans_step_5_recompute_centers.png}
    \end{figure}
\end{frame}
\begin{frame}{K-Means in action (cont.)}
    \begin{figure}
        \centering
        \includegraphics[scale=0.45]{pic/figs/kmeans_step_6_assignment.png}
    \end{figure}
\end{frame}
\begin{frame}{K-Means in action (cont.)}
    \begin{figure}
        \centering
        \includegraphics[scale=0.45]{pic/figs/kmeans_step_7_recompute_centers.png}
    \end{figure}
\end{frame}
\begin{frame}{K-Means in action (cont.)}
    \begin{figure}
        \centering
        \includegraphics[scale=0.45]{pic/figs/kmeans_step_8_assignment.png}
    \end{figure}
\end{frame}
\begin{frame}{K-Means in action (cont.)}
    \begin{figure}
        \centering
        \includegraphics[scale=0.45]{pic/figs/kmeans_step_9_recompute_centers.png}
    \end{figure}
\end{frame}
\begin{frame}{K-Means in action (cont.)}
    \begin{figure}
        \centering
        \includegraphics[scale=0.45]{pic/figs/kmeans_step_10_assignment.png}
    \end{figure}
\end{frame}
\begin{frame}{K-Means in action (cont.)}
    \begin{figure}
        \centering
        \includegraphics[scale=0.45]{pic/figs/kmeans_step_11_recompute_centers.png}
    \end{figure}
\end{frame}
\begin{frame}{K-Means in action (cont.)}
    \begin{figure}
        \centering
        \includegraphics[scale=0.45]{pic/figs/kmeans_step_12_assignment.png}
    \end{figure}
\end{frame}
\begin{frame}{K-Means in action (cont.)}
    \begin{figure}
        \centering
        \includegraphics[scale=0.45]{pic/figs/kmeans_step_13_recompute_centers.png}
    \end{figure}
\end{frame}
\begin{frame}{K-Means in action (cont.)}
    \begin{figure}
        \centering
        \includegraphics[scale=0.45]{pic/figs/kmeans_step_14_assignment.png}
    \end{figure}
\end{frame}


\begin{frame}{Convergence}
    \begin{itemize}
    \item How do we know K-Means will converge in a finite number of steps ?
        \item  In Assignment step we optimize \( J \) with respect to \( r_k(x) \).
        \item Because \( J \) is a linear function of \( r_k(x) \), since $x$s are not variables in training. this optimization can be performed easily to give a closed form solution.
        \item We need each $x$ to be at least in some cluster and terms involving different $x$ are independent, so for each $x$ we chose one of the the $k$ distance expressions that is the minimum. 
        \item  Now for $r_n$s fixed $J$ is a quadratic function of $\mu_k$ and by taking derivative we can miminze as:
$$
    
        \frac{\partial J}{\partial \mu_k} = 0 \implies 2 \sum_{i=1}^{N} r_k(x_i) \left(x_i - \mu_k \right) = 0
        
$$
$$

        \mu_k = \frac{\sum_{i=1}^{N}{r_k(x_i)x_i}}{\sum_{i=1}^{N} r_k(x_i)}
         
$$
        
    \end{itemize}
\end{frame}


\begin{frame}{Convergence (cont.)}
    \begin{itemize}
\item  Because each step reduces the value of the objective function J (or just does not increase it), convergence of the algorithm is assured.
\item  The convergence properties of the K-means algorithm were studied by MacQueen (1967).
\item We know each step will decrease the $J$ objective from its current value. Also there are a finite number of partitions, and $J$ is non-negative. So we must converge at some point, where the $J$ does not decrease anymore.
\item  Each Assignment and Updating step corresponds respectively to the E (expectation) and M (maximization) steps of the EM algorithm, which we will see later ???
    \end{itemize}

\end{frame}


\begin{frame}{K-Means convergence (cont.)}
    \begin{figure}
        \centering
        \includegraphics[scale=0.45]{pic/figs/distortion.png}
    \end{figure}
\end{frame}


\begin{frame}{Strengths}
    \begin{itemize}
        \item Simple: easy to understand and to implement.
        \item  Efficient: Time complexity: $O(tkn)$, 
where
\begin{itemize}
    
\item  $n$ is the number of data points, 
\item $k$ is the number of clusters, and 
\item $t$ is the number of iterations. 
\end{itemize}
– Since both k and t are usually small. k-means is considered a linear algorithm. (personally not sure about this ???)
        
    \end{itemize}
\end{frame}

\begin{frame}{Some Issues}
    \begin{itemize}
        \item k-Means always converges. What could go wrong ?
        \item K-means algorithm is a \textbf{heuristic}
        \item It requires initial centroids, and the choice is important. It could affect the $t$ in $O(tkn)$.
        \item the algorithm finds a local Minimum but it does not guarantee global minimum. This is highly affected by the initialization.
        \item Whats the solution ? some suggestions are:
        \begin{itemize}
            \item  variance-based split / merge
            \item Random centers from the data points with Multiple runs and select the best ones.
            \item initialization heuristics (k-means++ , Furthest Traversal)
            \item  Initializing with the suggested results of another method
        \end{itemize}
        
        
    \end{itemize}
\end{frame}

\begin{frame}{Some Issues (cont.)}
    \begin{itemize}
        \item In the begging we assumed $x_i \in \mathbb{R}^d$, which is not always the case. K-Means requires a space where sample \textbf{mean} is defined.
        \begin{itemize}
            \item A simple case is when we have categorical data. 
            \item  A suggested solution: k-mode - the centroid is represented by most frequent values.
        \end{itemize}
        \item The K is predefined.  
        \item The algorithm is sensitive to outliers
        \begin{itemize}
            \item Outliers are data points that are very far away from other data points. 
\item Outliers could be errors in the data recording or some special data points with very different values.
        \end{itemize}

    \end{itemize}
\end{frame}

\begin{frame}{Local optimum}
    \begin{itemize}
        \begin{figure}
            \centering
            \includegraphics[scale=0.55]{pic/original_data.png}
        \end{figure}
    \end{itemize}
\end{frame}
\begin{frame}{Local optimum (cont.)}
    \begin{minipage}{0.5\textwidth}
        \begin{figure}
            \centering
            \includegraphics[scale=0.55]{pic/optimal_clustering.png}
        \end{figure}
        \vfill
        \begin{center}
            Optimal clustering
        \end{center}
    \end{minipage}%
    \begin{minipage}{0.5\textwidth}
        \begin{figure}
            \centering
            \includegraphics[scale=0.55]{pic/possible_clustering.png}
        \end{figure}
        \vfill
        \begin{center}
            Possible clustering
        \end{center}
    \end{minipage}
\end{frame}


\begin{frame}{How many clusters?}
    \begin{figure}
        \centering
        \includegraphics[scale=0.5]{pic/how_many_clusters.png}
    \end{figure}
    \begin{tikzpicture}[remember picture,overlay]
        \node[anchor=south west, xshift=0.1cm, yshift=0.22cm] at (current page.south west) {
            \scriptsize Figure adapted from slides of Dr. Soleymani, Modern Information Retrieval Course, Sharif University of technology.
        };
    \end{tikzpicture}
\end{frame}

\begin{frame}{How many clusters? (cont.)}
    \begin{itemize}
        \item Number of clusters is usually given in advance in the problem of clustering. However; finding the \textbf{right} number of clusters is also a problem.
        \item There is a tradeoff between having better focus within each cluster or having too many clusters. We probably dont want each data to be a one-element cluster.
        \item \textbf{Optimization problem:} penalize having too much clusters
        \begin{itemize}
            \item Application dependent
        \end{itemize}
        \[
        K^* = \ arg \ min_k \  J(k) + \lambda k 
        \]
    \end{itemize}
\end{frame}

% \begin{frame}{others}
%     How do you find optimal K? (Is the elbow informative?)
% Outliers and pointless clusters and introducing k-means++
% Unknown number of clusters (regularization)
% External clustering criteria (purity, r-index, NMI, F-measure)

% \end{frame}

\begin{frame}{External criteria}
    External clustering criteria (purity, r-index, NMI, F-measure)
\end{frame}


\begin{frame}{Deal Breaker}
    \begin{itemize}
        \item Perhaps the most important problems is how k-means defines clusters. data might not be clustered in centered groups.
        \item As commonly seen in ML, an idea is to map this data into another vector space where thay can be clustered by our simple measures of similarity.
\item So lets take a look at clustering.

PICTURE

        
    \end{itemize}
\end{frame}


\section{Clustering}

\begin{frame}{Clustering}
    \begin{minipage}{0.95\textwidth}
        \begin{itemize}
        \item Assume we have a set of unlabeled data points $ \{ \mathbf{x}^{(i)} \}_{i=1}^N$.
        \item We intend to find \textbf{groups of similar objects} with respect to our need.
        \begin{itemize}
            \item For example all data points having most similar number of buys in a market.
        \end{itemize}
        \item It helps us to gain insight into structure of data prior to class design.
        \item Clustering could also help to compress and reduce data.
    \end{itemize}
    \end{minipage}%
    % \begin{minipage}{0.45\textwidth}
    %     \begin{itemize}
    %     \item Assume we have a set of unlabeled data points $ \{ \mathbf{x}^{(i)} \}_{i=1}^N$.
    %     \item We intend to find \textbf{groups of similar objects} with respect to our need.
    % \end{itemize}
    % \end{minipage}
\end{frame}

\begin{frame}{Clustering (cont.)}
    \begin{minipage}{0.55\textwidth}
        From another point of view, clusters are regions of high density that are separated from one another with regions of low density.
    % here we should talk what is good cluster ? should we even try to cluster based on density function ? should we focus of pure samples ? showcase the problem in which clustering could be both circle and a line extending of pmf (or pdf)
    \end{minipage}%
    \begin{minipage}{0.4\textwidth}
        \begin{figure}
            \centering
            \includegraphics[scale=0.5]{pic/clusteriong_pov.png}
        \end{figure}
    \end{minipage}
    \begin{tikzpicture}[remember picture,overlay]
        \node[anchor=south west, xshift=0.1cm, yshift=0.22cm] at (current page.south west) {
            \scriptsize Figure adapted from slides of Dr. Soleymani, Machine Learning course, Sharif University of technology.
        };
    \end{tikzpicture}
\end{frame}

\begin{frame}{Hard clustering vs Soft clustering}
    \begin{minipage}{0.55\textwidth}
        \begin{itemize}
        \item \textbf{Hard Clustering:} Each data point belongs to exactly one cluster
        \begin{itemize}
            \item more common and easier to do
        \end{itemize}
        \item \textbf{Soft Clustering}
    \end{itemize}
    \end{minipage}%
    \begin{minipage}{0.40\textwidth}
        \begin{figure}
            \centering
            \includegraphics[scale=0.5]{pic/hard_clustering.png}
        \end{figure}
    \end{minipage}
    \begin{tikzpicture}[remember picture,overlay]
        \node[anchor=south west, xshift=0.1cm, yshift=0.22cm] at (current page.south west) {
            \scriptsize Figure adapted from Machine Learning and Pattern Recognition, Bishop
        };
    \end{tikzpicture}  
\end{frame}

\begin{frame}{Hard clustering vs Soft clustering (cont.)}
    \begin{minipage}{0.55\textwidth}
        \begin{itemize}
        \item \textbf{Hard Clustering}
        \item \textbf{Soft Clustering:} Each data point can belong to multiple clusters.
        \begin{itemize}
            \item data point belongs to each cluster with a probability
        \end{itemize}
        \item \textbf{From now on, we will focus on problem of hard clustering}
    \end{itemize}
    \end{minipage}%
    \begin{minipage}{0.40\textwidth}
        \begin{figure}
            \centering
            \includegraphics[scale=0.5]{pic/soft_clustering.png}
        \end{figure}
    \end{minipage}
    \begin{tikzpicture}[remember picture,overlay]
        \node[anchor=south west, xshift=0.1cm, yshift=0.22cm] at (current page.south west) {
            \scriptsize Figure adapted from Machine Learning and Pattern Recognition, Bishop
        };
    \end{tikzpicture}  
\end{frame}

\begin{frame}{Hard clustering problem}
    \begin{itemize}
        \item We need to \textbf{partition} \( N \) data points to \( K \) clusters.
        \item Good flat clustering should have two important factors:
        \begin{itemize}
            \item Data points in a cluster should be similar (\textbf{High intra-cluster similarity})
            \item Data points in different clusters should be less similar (\textbf{Low inter-cluster similarity})
        \end{itemize}
        \item As mentioned earlier, each partitional clustering method requires a \textbf{similarity metric} between data points.
    \end{itemize}
\end{frame}

\begin{frame}{Similarity measure and Distance measure}
    \begin{minipage}{0.6\textwidth}
        \begin{itemize}
            \item Similarity measures are used to distinguish between similar and non-similar data points.
            \item We usually define similarity of two data point as \textbf{inverse of distance} between them.
            \item Using this definition, {hard clustering aims to put data points with less distance in same cluster.}
        \end{itemize}
    \end{minipage}%
    \begin{minipage}{0.35}
        \begin{figure}
            \centering
            \includegraphics[scale=0.4]{pic/Distance-based-clustering-approach.png}
        \end{figure}
    \end{minipage}
\end{frame}

\begin{frame}{Common similarity and distance measures}
    \begin{itemize}
        \item Assume \( p \) and \( q \) are two data points from $ \mathbb{R}^D$. most common similarity and distance measures in the problem of clustering are as follows:
        \begin{itemize}
            \item \textbf{Euclidean distance:} Most common measure of distance between two vectors doesn't matter.
            \[
            d^2(p,q)=\sqrt{{\sum_{i=1}^{D}{(p_i-q_i)^2}}}
            \]
            \item \textbf{Cosine similarity:} Most common measure of similarity when the magnitude of vectors does not change the similarity
            \[
            \textbf{similarity(\(p\),\(q\))}=\frac{p^T q}{||p|| \ . \ ||q||}
            \]
            \item \textbf{Manhattan distance:} Most common measure of distance when dimensions are not equally important
            \[
            d^2(p,q)=\sqrt{\sum_{i=1}^{D} |p_i-q_i|}
            \]
        \end{itemize}
    \end{itemize}
\end{frame}

\section{References}

\begin{frame}{Contributions}
\begin{itemize}
\item \textbf{This slide has been prepared thanks to:}
\begin{itemize}
    % \setlength{\itemsep}{10pt} % Adjust the value to control the spacing
    % \item \href{https://github.com/Mahan-Bayhaghi}{Mahan Bayhaghi}
\end{itemize}
\end{itemize}

\end{frame}

\begin{frame}[allowframebreaks]
    \bibliography{ref}
    \bibliographystyle{ieeetr}
    \nocite{*}
\end{frame}

\end{document}